\documentclass[11pt,a4paper,sans]{moderncv} 
\moderncvstyle{banking} 
\moderncvcolor{black} 
% Add this to make bold more visible without changing style:
\usepackage{xcolor}

% Subtle enhancement - makes bold text slightly darker
\let\oldtextbf\textbf
\renewcommand{\textbf}[1]{\oldtextbf{\color{black!90}#1}}

% Optional: highlight technical terms with slight emphasis
\newcommand{\tech}[1]{\textbf{#1}}
\usepackage[utf8]{inputenc}

% Adjusted scale slightly to ensure fit
\usepackage[scale=0.94, top=0.5cm, bottom=0.5cm]{geometry} 
\usepackage{fancyhdr}


% Bullet point spacing - Optimized
\usepackage{enumitem}
\setlist[itemize]{label=\textbullet, itemsep=0.15em, topsep=0.2em, leftmargin=3ex}

% HYPERLINK SETUP: Links are clickable but appear BLACK (Professional)
\usepackage{hyperref} 
\hypersetup{
    colorlinks=true,
    linkcolor=black,
    urlcolor=black,
    pdftitle={Amin Pakzad - Research Scientist CV}
}
% Add to preamble
\renewcommand*{\section}[1]{%
  \vspace{0.5cm}%
  {\Large\bfseries #1}%
  \par\nobreak
  \vspace{0.2cm}%
  \hrule height 1.2pt
  \vspace{1pt}
  \hrule height 0.4pt
  \par\nobreak
  \vspace{0.3cm}%
}
%--------------------------
% Personal data
%--------------------------
\name{Amin}{Pakzad}

\email{amin.pakzad.work@gmail.com}
\social[linkedin]{amin-pakzad-8a05b0127}
\social[github]{amnp95}
\extrainfo{Applied Scientist (ML) | Large-Scale Systems \& Model Development}

\begin{document}
\makecvtitle
\vspace*{-1.0cm}

%--------------------------
% Profile
%--------------------------
\section{Profile}
\cvitem{}{
\small{
Applied Scientist with a focus on massively parallel systems and distributed software architecture. Expert in scaling Python/C++ applications across thousands of cores and designing end-to-end pipelines for large-scale modeling. Interested in the intersection of high-performance engineering and applied machine learning
}}

%--------------------------
% Experience
%--------------------------
\section{Professional Experience}

\cventry{Jul 2023 -- Present}{Research Engineer / Applied Scientist}{SimCenter \& DesignSafe (NSF)}{Berkeley, CA}{}{
\begin{itemize}
  \item Designed and implemented end-to-end computational frameworks combining Python and C++ to execute large-scale workloads on distributed infrastructure (TACC clusters).
  \item Developed scalable data ingestion, preprocessing, and post-processing workflows to support high-throughput experimentation across thousands of parallel jobs.
  \item Optimized performance-critical components using MPI, enabling experiments to scale efficiently across thousands of CPU cores.
  \item Built robust execution pipelines with standardized configuration and result aggregation to support comparative model and simulation evaluation.
  \item Authored developer-facing documentation and Jupyter-based analysis notebooks to enable rapid experimentation by external users.
\end{itemize}}

\cventry{Jan 2022 -- Present}{Graduate Research Assistant / Core Developer}{University of Washington}{Seattle, WA}{}{
\begin{itemize}
  \item Architected \textbf{Femora}, a Python framework for automated model generation and experimentation, used as part of downstream ML and simulation pipelines.
  \item Implemented algorithms for large-scale geometry and data processing, reducing preprocessing time from weeks to minutes.
  \item Extended performance-critical C++ components to support adaptive control logic and improved numerical stability in long-running experiments.
  \item Designed binary I/O pipelines (HDF5, VTK) to efficiently store and analyze large experimental outputs.
  \item Built validation pipelines comparing simulated outputs against experimental datasets, enabling quantitative evaluation of model accuracy and generalization.
\end{itemize}}

%--------------------------
% Technical Skills
%--------------------------
\section{Technical Skills}
\cvitem{Machine Learning}{PyTorch, Scikit-learn, PINNs, model training \& evaluation, experiment design.}
\cvitem{Systems \& Performance}{MPI, OpenMP, CUDA, Linux, Docker, distributed computing, profiling.}
\cvitem{Languages}{Python (expert), C++ (advanced), MATLAB, SQL.}
\cvitem{Data \& Tools}{NumPy, Pandas, HDF5, Jupyter, Git, CI workflows.}

%----------------------------------------------------------------------------------------
%    EDUCATION
%----------------------------------------------------------------------------------------
\section{Education}
\cventry{Exp. 2026}{PhD Civil Engineering (Computational)}{University of Washington}{GPA: 3.85}{}{\small \textit{Thesis: High-Fidelity Dynamic Analysis utilizing Massively Parallel Computing (HPC).}}

\cventry{Exp. 2027}{M.S. Computer Science (Machine Learning)}{Georgia Tech}{GPA: 4.0}{}{}

\cventry{2024}{M.S. Applied Mathematics}{University of Washington}{GPA: 4.0}{}{}

%--------------------------
% Publications
%--------------------------
\section{Selected Software \& Publications}
\cvitem{2024}{McKenna, F., ... \textbf{Pakzad, A.}, et al. NHERI-SimCenter/EE-UQ v4.1.0 (Zenodo).}
\cvitem{2024}{\textbf{Pakzad, A.}, Arduino, P. Workflow for high-fidelity dynamic analysis (Pan-American Conference).}
\cvitem{2024}{\textbf{Pakzad, A.}, Arduino, P. Assessment of constitutive models (Soil Dynamics \& Earthquake Engineering).}

\end{document}